\documentclass[french]{book}
\usepackage{teach}
\usepackage{verbatim}
\usepackage{amsmath,amsfonts,amssymb}
\usepackage{pgfplots}
\usepackage{mwe}
\usepackage{yhmath} 
\usepackage{pstricks,pst-plot,pst-text,pst-tree,pst-eucl}

\newcommand{\arc}[1]{\wideparen{#1}}
\RequirePackage{graphicx}
\RequirePackage{ifpdf}
 \ifpdf
   \DeclareGraphicsRule{*}{mps}{*}{}
 \fi

\newcommand{\sysd}[2]{%
       \left\{
       \begin{aligned}
          #1\\
          #2\\
       \end{aligned}
       \right.%
       }
  \usepackage{fancyhdr}

  \usepackage{amsmath}
  \usepackage{amssymb}
  \usepackage{amsfonts}
  \usepackage{amsthm}
  \usepackage{mathrsfs}

  \usepackage{array}
  \usepackage{multicol}
  \setlength{\multicolsep}{3pt}

  \usepackage{graphicx}
  \graphicspath{{Figures/}}
  \usepackage{pstricks,pst-plot,pst-text,pst-tree,pst-eucl}
  \usepackage[all]{xy}

  \usepackage{fancybox}
  \usepackage{color}

%  \usepackage[frenchb]{babel}
  \usepackage{hyperref}

\usepackage{subfig}
\pgfplotsset{compat=newest}
\usepackage[all]{xy}
%\usepackage{geometry}
%\geometry{top=1cm, bottom=1cm, left=2cm, right=2cm}
\usepackage[utf8]{inputenc}
\usepackage[french]{babel}
\redefinecolor{thm}{title}{bgcolor}{alizarin}
\redefinecolor{prop}{title}{bgcolor}{meatbrown}
\redefinecolor{defn}{title}{bgcolor}{olivine}
\definecolor{alizarin}{rgb}{0.82, 0.1, 0.26}
\definecolor{meatbrown}{rgb}{0.9, 0.72, 0.23}
\definecolor{olivine}{rgb}{0.6, 0.73, 0.45}
\usepackage{mathrsfs}
%%
\newcommand{\R}{\mathbb {R}}
%\newcommand{\C}{\mathds {C}}
\newcommand{\N}{\mathbb {N}}
\newcommand{\Z}{\mathbb {Z}}
\newcommand{\Q}{\mathbb {Q}}
\newcommand{\D}{\mathbb {D}}
%%
\newcommand{\se}{\geq}
\newcommand{\ie}{\leq}
\newcommand{\tm}{\times}
%\newcommand{\ell}{l}
%\newcommand{\ell'}{l'}
%%
\newcommand{\barre}[1]{\overline{#1}}
\newcommand{\ds}{\displaystyle}
\newcommand{\limn}{\lim_{n\rightarrow +\infty}}
\newcommand{\limo}[1][x]{\ds\lim_{#1\rightarrow 0}}
\newcommand{\limite}[2][x]{\ds\lim_{#1\rightarrow #2}}
\newcommand{\limpinf}[1][x]{\ds\lim_{#1\rightarrow +\infty}}
\newcommand{\limminf}[1][x]{\ds\lim_{#1\rightarrow -\infty}}
\newcommand{\limiteg}[2][x]{\ds\lim_{#1\xrightarrow{<}#2}}
\newcommand{\limited}[2][x]{\ds\lim_{#1\xrightarrow{>}#2}}
%%
\newcommand{\vect}[1]{\overrightarrow{#1}}
\newcommand{\oij}{(\textit{O},{\overrightarrow{i}},{\overrightarrow{j}})}
%
\newcommand{\cp}[2]{%
        \begin{pmatrix}
        #1\\
        #2
        \end{pmatrix}%
        }
%
\newcommand{\calb}{\mathscr{B}}
\newcommand{\calc}{\mathscr{C}}
\newcommand{\cald}{\mathscr{D}}
\newcommand{\cale}{\mathscr{E}}
\newcommand{\calf}{\mathscr{F}}
\newcommand{\calp}{\mathscr{P}}
\newcommand{\cals}{\mathscr{S}}
\newcommand{\calt}{\mathscr{T}}
%
\newcommand{\suite}[1][u]{\left( #1_{n} \right)_{n\in\N}}
\newcommand{\suitep}[1][u]{\left( #1_{n} \right)_{n\in\N^{*}}}
\usetikzlibrary{arrows}

\newcommand{\poly}[1][a]{#1_0+#1_1x+\cdots+#1_nx^n}
\usepackage{tabularx}
\usepackage{array}
\usepackage{pst-tree}
\usepackage{pifont}
%\usepackage[cp1252]{inputenc}



\newcommand{\V}{\overrightarrow}
\newcommand{\Rep}{(O;\V{i};\V{j})}
%\newcommand{\Int}[2]{[\, #1 ; #2 \,]}
%\newcommand{\Reel}{\mathfrak{Re}} 
%\newcommand{\Ima}{\mathfrak{Im}} 
%\renewcommand{\Abs}[1]{\left \lvert#1\right \rvert}
\newcommand{\Cnp}{\begin{pmatrix} n\\p \end{pmatrix}}
\newcommand{\C}[2]{\begin{pmatrix} #1\\#2 \end{pmatrix}}
\newcommand{\Syst}[2]{\left\{\begin{array}{lcccc} #1 \\ #2 \end{array}\right.}
\renewcommand{\arraystretch}{1.2} 
\newcolumntype{C}[1]{>{\centering\arraybackslash}p{#1cm}}


\begin{document}


\tableofcontents










%%%%%%%%%%%%%%%%%%%%%%%%%%%%%%%%%%%%%%%%%%%%%%%%%%%%%
%%%   I   %%%%%%%%%%%%%%%%%%%%%%%%%%%%%%%%%%%%%%%%%%%
%%%%%%%%%%%%%%%%%%%%%%%%%%%%%%%%%%%%%%%%%%%%%%%%%%%%%
\chapter{Statistiques à deux variables}
Le problème qui se pose dans les séries statistiques à deux variables est principalement celui du lien qui existe ou non entre chacune des variables. 
\section{Position du problème. Vocabulaire}

Par soucis de clarté, ce cours est élaboré à partir de l'exemple suivant :



\begin{exemple}
Le tableau suivant donne l'évolution du nombre d'adhérents d'un club de rugby de $2001$ à $2006$.
\begin{center}
   \begin{tabular}{|c|*{6}{C{1}|}}
      \hline
      Année & $2001$ & $2002$ & $2003$ & $2004$ & $2005$ & $2006$ \\
      \hline
      Rang $x_{i}$ & $1$ & $2$ & $3$ & $4$ & $5$ & $6$ \\
      \hline
      Nombre d'adhérents $y_{i}$ & $70$ & $90$ & $115$ & $140$ & $170$ & $220$ \\
      \hline
   \end{tabular}
\end{center}
Le but est d'étudier cette série statistique à deux variables (le rang et le nombre d'adhérents) afin de prévoir l'évolution du nombre d'adhérents pour les années suivantes.
\end{exemple}

%%%   I.1   %%%%%%%%%%%%%%%%%%%%%%%%%%%%%%%%%%%%%%%%%%%
\subsection{Nuage de points}

La première étape consiste à réaliser un graphique qui traduise les deux séries statistiques ci-dessus.

\begin{defn}
   Soit $X$ et $Y$ deux variables statistiques numériques observées sur $n$ individus.\\
   Dans un repère orthogonal $\Rep$, l'ensemble des $n$ points de coordonnées $(x_i,y_i)$ forme le \underline{nuage de points} associé à cette série statistique.
\end{defn}

Dans notre exemple, si on place le rang en abscisses, et le nombre d'adhérents en ordonnées, on peut représenter par un point chaque valeur. On obtient ainsi une succession de points, dont les coordonnées sont $(1;70)$, $(2;90)$, ... $(6;220)$, forment un \underline{nuage de points}.



\begin{exemple}
   Dans le plan muni d'un repère orthogonal d'unités graphiques : $2$~cm pour une année sur l'axe des abscisses et $1$~cm pour $20$~adhérents sur l'axe des ordonnées, représenter le nuage de points associé à la série $(x_{i};y_{i})$.
\end{exemple}

\begin{color}{blue}
   \begin{multicols}{2}
   T.I.
   \begin{itemize}
      \item[$\bullet$] Touche \fbox{STAT}
      \item[$\bullet$] Menu \fbox{EDIT}
      \item[$\bullet$] Entrer les valeurs $x_i$ dans $L_1$
      \item[$\bullet$] Entrer les valeurs $y_i$ dans $L_2$
      \item[$\bullet$] Règler les valeurs du repère avec la touche \fbox{WINDOWS}
      \item[$\bullet$] Appuyer sur la touche \fbox{TRACE}
   \end{itemize}
   Casio
   \begin{itemize}
      \item[$\bullet$] Menu \fbox{STAT}
      \item[$\bullet$] Entrer les valeurs $x_i$ dans $List 1$
      \item[$\bullet$] Entrer les valeurs $y_i$ dans $List 2$
      \item[$\bullet$] Choisir \fbox{GRPH}
      \item[$\bullet$] Règler les paramètres avec \fbox{SET}
      \item[$\bullet$] Choisir \fbox{GPH1}
   \end{itemize}
   \end{multicols}
\end{color}

\begin{center}
   \psset{unit=0.9}
   \begin{pspicture}(-1,-0.5)(18,14.5)
      \psset{xunit=2,yunit=0.05}
      \psgrid[subgriddiv=1,griddots=2,gridlabels=0](0,0)(8,280)
      \psaxes[Dy=20]{->}(0,0)(8.2,282)
      \psline[linestyle=dotted](0,20)(8,20)
      \psline[linestyle=dotted](0,40)(8,40)
      \psline[linestyle=dotted](0,60)(8,60)
      \psline[linestyle=dotted](0,80)(8,80)
      \psline[linestyle=dotted](0,100)(8,100)
      \psline[linestyle=dotted](0,120)(8,120)
      \psline[linestyle=dotted](0,140)(8,140)
      \psline[linestyle=dotted](0,160)(8,160)
      \psline[linestyle=dotted](0,180)(8,180)
      \psline[linestyle=dotted](0,200)(8,200)
      \psline[linestyle=dotted](0,220)(8,220)
      \psline[linestyle=dotted](0,240)(8,240)
      \psline[linestyle=dotted](0,260)(8,260)
      \psdots(1,70)(2,90)(3,115)(4,140)(5,170)(6,220) 
      \psdots[dotstyle=x](2,91.7)(5,176.7)(3.5,134.2)            
      \rput(2,100){$G_1$}
      \rput(3.5,143){$G$}
      \rput(5,187){$G_2$}
      \psline[linecolor=green](0,35.1)(8,261.5)
      \rput(7.2,230){\textcolor{green}{$\mathcal{D}_1$}} 
      \psdots[dotstyle=+](0,32.7)(8,264.7)
      \psline[linecolor=red](0,32.7)(8,264.7)
      \rput(7.5,260){\textcolor{red}{$\mathcal{D}_2$}} 
      \pscurve[linecolor=blue](0,57.1)(1,71.4)(2,89.4)(3,111.8)(4,139.9)(5,175)(6,218.9)(7,273.9)(7.1,280)
      \rput(6.5,260){\textcolor{blue}{$\mathcal{C}_f$}} 
      \rput(8.4,-8){Rang}
      \rput(0.3,288){Nombre d'adhérents}
   \end{pspicture}
\end{center}


%%%   I.2  %%%%%%%%%%%%%%%%%%%%%%%%%%%%%%%%%%%%%%%%%%%
\subsection{Le problème de l'ajustement}

Le nuage de points associé à une série statistique à deux variables donne donc immédiatement des informations de nature qualitatives.\\
Pour en tirer des informations plus quantitatives, il nous faut poser le problème de l'ajustement.

Le tracé met en évidence la possibilité de "reconnaître" graphiquement la possibilité d'une relation fonctionnelle entre les deux grandeurs observées (ici rang et nombre d'adhérent).\\
Le problème de l'établissement d'une relation fonctionnelle entre les deux séries est le \underline{problème de l'ajustement}. 

%%%   I.3   %%%%%%%%%%%%%%%%%%%%%%%%%%%%%%%%%%%%%%%%%%%
\subsection{Point moyen}

\begin{defn}
   Soit une série statistique à deux variables, $X$ et $Y$, dont les valeurs sont des couples $(x_i;y_i)$.\\
   On appelle \underline{point moyen} de la série le point $G$ de coordonnées
   \begin{itemize}
      \item[$\bullet$] $x_G=\dfrac{x_1+x_2+ \dots + x_n}{n}$.
      \item[$\bullet$] $y_G=\dfrac{y_1+y_2+ \dots + y_n}{n}$.
   \end{itemize}
\end{defn}

\begin{exemple}
   Déterminer les coordonnées des points moyens suivants :
   \begin{itemize}
      \item[$\bullet$] $G_1$ des années allant de $2001$ à $2003$,
      \item[$\bullet$] $G_2$ des années allant de $2004$ à $2006$,
      \item[$\bullet$] $G$, point moyen du nuage de points tout entier.
   \end{itemize}
\end{exemple}

Calcul des coordonnées de $G_1$ : $\Syst{x_{G_1} & = & \frac{1+2+3}{3} & = & 2}{y_{G_1} & = & \frac{70+90+115}{3} & = & 91,7}$ \qquad donc, ~ \fbox{$G_1(~2~;~91,7~)$}.

Calcul des coordonnées de $G_2$ : $\Syst{x_{G_2} & = & \frac{4+5+6}{3} & = & 5}{y_{G_2} & = & \frac{140+170+220}{3} & = & 176,7}$ \qquad donc, ~ \fbox{$G_2(~5~;~176,7~)$}.

Calcul des coordonnées de $G$ : $\Syst{x_G & = & \frac{1+2+3+4+5+6}{6} & = & 3,5}{y_G & = & \frac{70+90+115+140+170+220}{3} & = & 134,2}$ \qquad donc, ~ \fbox{$G(~3,5~;~134,2~)$}.


%%%%%%%%%%%%%%%%%%%%%%%%%%%%%%%%%%%%%%%%%%%%%%%%%%%%%
%%%   II   %%%%%%%%%%%%%%%%%%%%%%%%%%%%%%%%%%%%%%%%%%
%%%%%%%%%%%%%%%%%%%%%%%%%%%%%%%%%%%%%%%%%%%%%%%%%%%%%
\section{Ajustements}


%%%   II.1   %%%%%%%%%%%%%%%%%%%%%%%%%%%%%%%%%%%%%%%%%%%
\subsection{Ajustement à la règle}

On se propose, à partir des résultats obtenus, de faire des prévisions pour les années à venir. \\
Un poyen d'y parvenir est de tracer au juger une droite $\mathcal{D}$ passant le plus près possible des points du nuage et d'en trouver l'équation du type $y=ax+b$.


%%%   II.2   %%%%%%%%%%%%%%%%%%%%%%%%%%%%%%%%%%%%%%%%%%
\subsection{Méthode de Mayer}

Cet ajustement consiste à déterminer la droite passant par deux points moyens du nuage de point.


\begin{exemple}
   Déterminer l'équation de la droite $\mathcal{D}_1$ qui passe par les points moyens $G_1$ et $G_2$ et la tracer sur le graphique précédent.
\end{exemple}

La droite $\mathcal{D}_1$ n'est pas parallèle à l'axe des ordonnées, elle a donc pour équation $y=ax+b$ avec : \\
$a=\dfrac{y_{G_2}-y_{G_1}}{x_{G_2}-x_{G_2}}=\frac{176,7-91,7}{5-2}=28,3$.

De plus, elle passe par le point $G_1(~2~;~91,7~)$ d'où : \\
$y_{G_1}=ax_{G_1}+b \quad \Rightarrow \quad 91,7=28,3\times 2+b \quad \Rightarrow \quad b=35,1$.

Conclusion : \fbox{$\mathcal{D}_1 : y=28,3x+35,1$}.

Pour tracer $\mathcal{D}_1$, il suffit de placer $G_1$ et $G_2$ puis de tracer la droite qui les relie.


%%%   II.3   %%%%%%%%%%%%%%%%%%%%%%%%%%%%%%%%%%%%%%%%%%
\subsection{Méthode des moindres carrés}

Il s'agit d'obtenir une droite équidistante des points situés de part et d'autre d'elle-même. \\
Pour réaliser ceci, on cherche à minimiser la somme des distances des points à la droite au carré.

On considère une série statistique à deux variables représentée par un
nuage justifiant un ajustement affine.



\begin{defn}
   Dans le plan muni d'un repère orthonormal, on considère un nuage de $n$ points de coordonnées $(x_i;y_i)$.\\
   La droite $\mathcal{D}$ d'équation $y=ax+b$ est appelée \underline{droite de régression} de $y$ en $x$ de la série statistique ssi la quantité suivante est minimale :
   $$\sum_{i=1}^n~(M_iQ_i)^2=\sum_{i=1}^n [y_i-(ax_i+b)]^2$$
\end{defn}



\begin{center}
   \psset{unit=0.8}
   \begin{pspicture}(-0.5,-1)(10,6.5)
      \psaxes[labels=none,ticks=none](0,0)(-0.5,-0,5)(10,7)
      %\psgrid[subgriddiv=1,griddots=5,gridlabels=0](0,0)(10,7)
      \psdots(1,2.5)(3,2)(5,5.5)(7,4.5)(9,7)
      \psline[linecolor=red](0,1)(9,7)
      \psline[linestyle=dashed](1,1.7)(1,2.5)
      \psline[linestyle=dashed](3,3)(3,2)
      \psline[linestyle=dashed](5,4.3)(5,5.5)
      \psline[linestyle=dashed](7,5.7)(7,4.5)
      \psline[linestyle=dotted](0,5.5)(5,5.5)(5,0)
      \psline[linestyle=dotted](0,4.3)(5,4.3)
      \psdot[linecolor=red](5,4.3)
      \rput(-1,4.3){$ax_i+b$} 
      \rput(-0.5,5.5){$y_i$}
      \rput(5.5,5.8){$M_i$}
      \rput(5.5,4){\textcolor{red}{$Q_i$}}
      \rput(8.5,6){\textcolor{red}{$\mathcal{D}$}}
      \rput(5,-0.5){$x_i$}   
   \end{pspicture}
\end{center}

\begin{rem}
   Il serait tout aussi judicieux de s'intéresser à la droite $\mathcal{D}'$ qui minimise la quantité $\displaystyle{\sum_{i=1}^n [x_i-(ay_i+b)]^2}$. \\
   Cette droite est appelée droite de régression de $x$ en $y$.
\end{rem}



\begin{defn}
   On appelle \underline{covariance} de la série statistique double de variables $x$ et $y$ le nombre réel
$$cov(x,y)=\sigma_{xy}=\frac1n\sum_{i=1}^n (x_i - \bar x) (y_i - \bar y).$$
\end{defn}

Pour les calculs, on pourra aussi utiliser :
$$\sigma_{xy}=\frac1n\sum_{i=1}^n x_i y_i-\bar x\bar y.$$

\begin{rem}
   On a : \quad $cov(x,x)=\sigma_{x^2}=V(x)=[\sigma(x)]^2$.
\end{rem}



\begin{prop}
   La droite de régression $D$ de $y$ en $x$ a pour équation $y=ax+b$ où 
   
  %$$\Syst{a=\dfrac{\sigma_{xy}}{\sigma_{x}^2}}{b\rm{~~vérifie~~}\bar y=a\bar x+b.}$$
  
  $$\Syst{a=\dfrac{\sigma_{xy}}{\sigma_{x}^2}}{b \; \text{vérifie} \; \bar y = a\bar x +b}$$
\end{prop}

\begin{rem}
   Les réels $a$ et $b$ sont donnés par la calculatrice.
\end{rem}

\begin{color}{blue}
\begin{multicols}{2}
   T.I.
   \begin{itemize}
      \item[$\bullet$] Touche \fbox{STAT}
      \item[$\bullet$] Menu \fbox{CALC}
      \item[$\bullet$] Item \fbox{LinReg}
      \item[$\bullet$] LinReg $L_1,L_2$ \\
   \end{itemize}
   Casio
   \begin{itemize}
      \item[$\bullet$] Menu \fbox{STAT}
      \item[$\bullet$] Item \fbox{CALC}
      \item[$\bullet$] Règler les paramètres avec\fbox{set}
      \item[$\bullet$] Item \fbox{REG}
      \item[$\bullet$] Choisir \fbox{X}
   \end{itemize}
\end{multicols}
\end{color}



\begin{prop}
   Le point moyen $G$ du nuage appartient toujours à la droite de régression de $y$ en $x$.
\end{prop}



\begin{exemple}
   Déterminer une équation de la droite d'ajustement $\mathcal{D}_2$ de $y$ en $x$ obtenue par la méthode des moindres carrés et la tracer sur le graphique précédent.
\end{exemple}

La calculatrice donne $\mathcal{D}_2:y=ax+b$ avec $a=29$ et $b=32,7$.

Conclusion : \fbox{$\mathcal{D}_2:y=29x+32,7$}

Pour tracer la droite $\mathcal{D}_2$, il faut choisir deux points (au moins) sur cette droite.\\
Par exemple : \quad 
\begin{tabular}{C{1}|C{1}|C{1}}
   $x$ & $0$ & $8$ \\
   \hline
   $y$ & $32,7$ & $264,7$ \\
\end{tabular}
\quad , les placer dans le repère puis tracer la droite.

%%%   II.4   %%%%%%%%%%%%%%%%%%%%%%%%%%%%%%%%%%%%%%%%%%
\subsection{Ajustement exponentiel}%%%%%%%%%%%%%%%%%%%%%

On remarque qu'un ajustement affine ne semble pas très approprié pour ce nuage de points à partir de $2006$, on se propose de déterminer un ajustement plus juste.



\begin{exemple}
   On pose $z=\ln y$. Recopier et compléter le tableau suivant en arrondissant les valeurs de $z_{i}$ au millième. 
   \begin{center}
      \begin{tabular}{|*{7}{C{1}|}}
         \hline
         $x_{i}$ & $1$ & $2$ & $3$ & $4$ & $5$ & $6$ \\
         \hline
         $z_{i}$ & $4,248$ & & & & & \\
         \hline
      \end{tabular}
   \end{center}
\end{exemple}

Il suffit de calculer $\ln y_i$ pour chaque caleur de $i$ :
\begin{center}
   \begin{tabular}{|*{7}{C{1}|}}
      \hline
      $x_{i}$ & $1$ & $2$ & $3$ & $4$ & $5$ & $6$ \\
      \hline
      $z_{i}$ & $4,248$ & $4,500$ & $4,745$ & $4,942$ & $5,136$ & $5,394$ \\
      \hline
   \end{tabular}
\end{center}

On peut déterminer les éléments de ce tableau grâce à la calculatrice :

\begin{color}{blue}
\begin{multicols}{2}
   T.I.
   \begin{itemize}
      \item[$\bullet$] Touche \fbox{STAT}
      \item[$\bullet$] Menu \fbox{EDIT}
      \item[$\bullet$] Se placer dans \fbox{$L_3$}
      \item[$\bullet$] Entrer la formule "$=\ln L_2$"
   \end{itemize}
   Casio
   \begin{itemize}
      \item[$\bullet$] Touche \fbox{STAT}
      \item[$\bullet$] Menu \fbox{EDIT}
      \item[$\bullet$] Se placer dans \fbox{$List3$}
      \item[$\bullet$] Entrer la formule "$=\ln List2$"
   \end{itemize}
\end{multicols}
\end{color}




\begin{exemple}
   Déterminer une équation de la droite d'ajustement $\mathcal{D}_3$ de $z$ en $x$ obtenue par la méthode des moindres carrés.  
\end{exemple}

La manipulation à la calculatrice est la même que précédemment, en oubliant pas de changer les paramètres.

La calculatrice donne $\mathcal{D}_3:z=ax+b$ avec $a=0,224$ et $b=4,045$.

Conclusion : \fbox{$\mathcal{D}_3:z=0,224x+4,045$}.



\begin{exemple}
   Dans ce cas, en déduire la relation qui lie $y$ à $x$ puis tracer la courbe représentative de la fonction $y=f(x)$.
\end{exemple}

On a $\Syst{z & = & 0,224x+4,045}{z & = & \ln y}$ \quad donc : \quad $\ln y=0,224x+4,045$

On compose par la fonction exponentielle :
\begin{tabular}[t]{ccl}
   $e^{\ln y} $& $=$ &$ e^{0,224x+4,045}$ \\
    &$ = $&$ (e^{0,224})^x\times e^{4,045}$ \\
    &$ =$ & $(1,251)^x\times 57,111$ \\
\end{tabular}

Conclusion : \fbox{$y=57,111\times1,251^x$}.

Pour tracer la courbe, il suffit de placer des points, par exemple grâce au tableau de valeurs de la calculatrice.


%%%   II.5   %%%%%%%%%%%%%%%%%%%%%%%%%%%%%%%%%%%%%%%%%%
\subsection{Comparaison}

Grâce aux trois derniers ajustements, on peut évaluer ce qui se passera plus tard, comparons les :



\begin{exemple}
   En supposant que les ajustements restent valables pour les années suivantes, donner une estimation du nombre d'adhérents en $2007$ suivant les trois méthodes.
\end{exemple}

Dans tous les cas, il faut calculer $y$ lorsque $x$ correspond à l'année $2007$, c'est à dire au rang $7$.

\begin{itemize}
   \item[$\bullet$] Méthode de Mayer : $y=28,3\times7+35,1=233,2$ soit environ \fbox{$233$ adhérents}.
   \item[$\bullet$] Ajustement affine : $y=29\times 7+32,7=235,7$ soit environ \fbox{$236$ adhérents}.
   \item[$\bullet$] Ajustement exponentiel : $y=57,112\times1,024^7=273,9$ soit environ \fbox{$274$ adhérents}.
\end{itemize}



\begin{exemple}
   En $2007$, il y a eu $280$ adhérents. Lequel des trois ajustements semble le plus pertinent ?
\end{exemple}

Le troisième ajustement semble le plus pertinent puisqu'il se rapporche le plus de la réalité.

%%%%%%%%%%%%%%%%%%%%%%%%%%%%%%%%%%%%%%%%%%%%%%%%%%%%%
%%%   III   %%%%%%%%%%%%%%%%%%%%%%%%%%%%%%%%%%%%%%%%%
%%%%%%%%%%%%%%%%%%%%%%%%%%%%%%%%%%%%%%%%%%%%%%%%%%%%%
\section{Coefficient de corrélation linéaire}

\begin{defn}
   Le \underline{coefficient de corrélation linéaire} d'une série statistique de variables $x$ et $y$ est le nombre $r$ défini par :
   $$r=\frac{\sigma_{xy}}{\sigma(x)\times\sigma(y)}.$$
\end{defn}

Ce coefficient sert à mesurer la qualité d'un ajustement affine.


\textbf{Interprétation graphique :} \\
Plus le coefficient de régression linéaire est proche de $1$ en valeur absolue, meilleur est l'ajustement linéaire. \\
Lorque $r=\pm1$, la droite de régression passe par tous les points du nuage, qui sont donc alignés.


\begin{exemple}
   Déterminer le coefficient de corrélation linéaire dans le cas de l'ajustement affine (entre $x$ et $y$), puis exponentiel (entre $x$ et $z$). Quel est l'ajustemet le plus juste ?
\end{exemple}

Grâce à la calculatrice, on trouve successivement $r_2=0,987$ puis $r_3=0,999$. \\
Ce qui est conforme à ce que nous avions déduit précédemment, à savoir que l'ajustement exponentiel est plus fiable pour ce cas.



\begin{prop}
   Le  coefficient de corrélation linéaire $r$ vérifie $-1\leq r\leq 1$.
\end{prop}

 
\end{document}